\chapter{z-Transformation von Signalen}\label{sec:z-transformation-von-signalen}

Die Laplace-Transformation wird zur Beschreibung zeitkontinuierlicher Signale verwendet. Sie eignet sich zur Lösung linearer Differentialgleichungen mit konstanten Koeffizienten und wird zur Interpretation von linearen, zeitinvarianten Systemen mit Übertragungsfunktionen eingesetzt. Die Interpretation der Übertragungsfunktionen führt zu Aussagen wie Stabilität, Schwingungsneigung, Kausalität und Sprungfähigkeit von Systemen. Die z-Transformation ist für zeitdiskrete Anwendungen das Pendant zur Laplace-Transformation. Sie wird in diesem Kapitel eingeführt und ist eine wichtige Voraussetzung für die Beschreibung zeitdiskreter Systeme und den Entwurf zeitdiskreter Filter.

Nach der Definition werden einige Korrespondenzen über die Definitionsgleichung bestimmt, und es wird ein Zusammenhang zwischen der z-Transformation und der Laplace-Transformation hergestellt.

Die eher aufwendige Bestimmung von Korrespondenzen über die Definitionsgleichung kann vermieden werden, wenn die vorliegende Funktion mit Rechenregeln auf Funktionen mit bekannten Korrespondenzen zurückgeführt werden kann. Die dazu notwendigen Rechenregeln werden hergeleitet und ihr Nutzen an Beispielen aufgezeigt.

Die bei technischen Anwendungen entstehenden z-Transformierten sind typischerweise gebrochen rationale Funktionen. Die Rücktransformation vom z-Bereich in den Zeitbereich kann grundsätzlich über ein Umkehrintegral erfolgen. Da dieser Weg aufwendig ist und Kenntnisse in der Funktionentheorie voraussetzt, wird er in der Praxis umgangen. Stattdessen wird die z-Transformierte in Partialbrüche zerlegt, die mit Hilfe der angesprochenen Rechenregeln und einiger Korrespondenzen in den Zeitbereich transformiert werden können.

\section{Grundlagen der z-Transformation}\label{sec:5.1}

\subsection{Definitionsgleichung der z-Transformation}\label{sec:definitionsgleichung-der-z-transformation}

Die z-Transformation ist das Äquivalent zur Laplace-Transformation. Zur Herleitung der Definitionsgleichung wird deshalb von der Laplace-Transformation ausgegangen. Wie bei der Laplace-Transformation werden nur kausale Signale betrachtet. Nach den Darstellungen zur Signalabtastung wird ein ideal abgetastetes kausales Signal über die Multiplikation des kontinuierlichen Signals mit einer Abtastfunktion dargestellt:
\begin{align}
x_{A}(t) &= \sum\limits_{k = 0}^{\infty}{x\left( k  T_{A} \right)  \delta\left( t - k  T_{A} \right)} \label{eq:5.1}
\end{align}
Bei \(x_{A}(t)\) handelt es sich um ein kontinuierliches Signal. Deshalb kann die Laplace-Transformierte berechnet werden zu
\begin{align}
X_{A}(s) &= L\left\{ \sum\limits_{k = 0}^{\infty}{x\left( k  T_{A} \right)  \delta\left( t - k  T_{A} \right)} \right\} \label{eq:5.2}
\end{align}
Da die Laplace-Transformation eine lineare Operation ist und der Term \(x_{A}\)(kT\textsubscript{A}) nicht von der Zeit t abhängt, kann die Laplace-Transformierte umgerechnet werden in
\begin{align}
{L\left\{ \sum\limits_{k = 0}^{\infty}{x(k  T_{A})  \delta(t - k  T_{A})} \right\} = \sum\limits_{k = 0}^{\infty}{x\left( k  T_{A} \right)  L\left\{ \delta\left( t - k  T_{A} \right) \right\}} = \sum\limits_{k = 0}^{\infty}{x\left( k  T_{A} \right)  e^{- k  T_{A}  s}}
}{			 = \sum\limits_{k = 0}^{\infty}{x[ k]  \left( e^{T_{A}  s} \right)^{- k}}} \label{eq:5.3}
\end{align}
Dabei ist die Variable s eine komplexe Zahl mit Real- und Imaginärteil. Durch die Substitution von
\begin{align}
z &= e^{T_{A}  s} \label{eq:5.4}
\end{align}
kann die Schreibweise vereinfacht werden zu
\begin{align}
X(z) &= \sum\limits_{k = 0}^{\infty}{x[ k]  z^{- k}} \label{eq:5.5}
\end{align}
Diese Gleichung ist die Definitionsgleichung der z-Transformation für kausale Signale. Ähnlich wie bei den vorangegangenen Transformationen wird für die z-Transformation die Schreibweise
\begin{align}
Z\left\{ x[ k] \right\} &= X(z) \label{eq:5.6}
\end{align}
und für die inverse z-Transformation die Schreibweise
\begin{align}
Z^{- 1}\left\{ X(z) \right\} &= x[ k] \label{eq:5.7}
\end{align}
eingeführt. Alternativ kann das Hantel-Symbol verwendet werden.
\begin{align}
x[ k] \origof X(z) \label{eq:5.8}
\end{align}
\subsection{z-Transformation grundlegender Signale}\label{sec:z-transformation-grundlegender-signale}

Wie bei der Laplace-Transformation werden zunächst einige Korrespondenzen der z-Transformation direkt über die Definitionsgleichung der z-Transformation berechnet. Dabei ergibt sich die z-Transformierte eines Signals durch Auswertung von endlichen und unendlichen Reihen.

\subsubsection{Diskrete Impulsfolge}\label{sec:diskrete-impulsfolge}

Aus der Definition der diskreten Impulsfolge
\begin{align}
x[ k] = \delta[ k] = \begin{dcases}
    1 & \text{f}\text{ü}\text{r k } = \text{ 0} \\
    0 & \text{f}\text{ü}\text{r ganzzahlige  k } \neq \text{ 0}
\end{dcases}\ \label{eq:5.9}
\end{align}
ergibt sich die z-Transformierte durch Einsetzen in die Definitionsgleichung zu
\begin{align}
X(z) &= \sum\limits_{k = 0}^{\infty}{\delta[ k]  z^{- k}} = 1  z^{0} = 1 \label{eq:5.10}
\end{align}
Ist der Impuls um \(k_{0}\) verschoben, ändert sich die z-Transformierte zu
\begin{align}
X(z) &= \sum\limits_{k = 0}^{\infty}{\delta\left[ k - k_{0} \right]  z^{- k}} = 1  z^{- k_{0}} = z^{- k_{0}} \label{eq:5.11}
\end{align}
Der Vergleich der z-Transformierten von Impuls und verschobenem Impuls ist ein erster Hinweis auf eine Verschiebungsregel.

\subsubsection{Diskrete Rechteckfolge}\label{sec:diskrete-rechteckfolge}

Für die Rechteckfolge mit der Definition
\begin{align}
x[ k] &= \sigma[ k] - \sigma[ k - K] \label{eq:5.12}
\end{align}
wird die Summe in der Definitionsgleichung der z-Transformation endlich. Es gilt
\begin{align}
X(z) &= \sum\limits_{k = 0}^{\infty}{\left( \sigma[ k] - \sigma[ k - K] \right)  z^{- k}} = \sum\limits_{k = 0}^{K - 1}z^{- k} \label{eq:5.13}
\end{align}
Die Auswertung der endlichen Summe mit den Rechenregeln für endliche geometrische Reihen
\begin{align}
\sum\limits_{k = 0}^{K - 1}q^{k} &= \frac{1 - q^{K}}{1 - q} \label{eq:5.14}
\end{align}
ergibt die z-Transformierte
\begin{align}
X(z) &= \sum\limits_{k = 0}^{\infty}{\left( \sigma[ k] - \sigma[ k - K] \right)  z^{- k}} = \frac{1 - z^{- K}}{1 - z^{- 1}} \label{eq:5.15}
\end{align}
\subsubsection{Diskrete Sprungfolge}\label{sec:diskrete-sprungfolge}

Die z-Transformation der Sprungfolge mit der Definition
\begin{align}
x[ k] = \sigma[ k] = \begin{dcases}
    1 & \text{f}\text{ü}\text{r ganzzahlige k } \geq \text{ 0} \\
    0 & \text{f}\text{ü}\text{r ganzzahlige }k < \text{ 0}
\end{dcases}\ \label{eq:5.16}
\end{align}
errechnet sich direkt durch Einsetzen in die Definitionsgleichung zu
\begin{align}
X(z) &= \sum\limits_{k = 0}^{\infty}{\sigma[ k]  z^{- k}} = \sum\limits_{k = 0}^{\infty}z^{- k} = \sum\limits_{k = 0}^{\infty}\left( z^{- 1} \right)^{k} \label{eq:5.17}
\end{align}
Die Berechnung der Reihe kann auf die unendliche geometrische Reihe zurückgeführt werden. Die geometrische Reihe ist definiert als:
\begin{align}
1 + q + q^{2} + ... &= \sum\limits_{k = 0}^{\infty}q^{k} \label{eq:5.18}
\end{align}
Die Konvergenz der Reihe ist von dem Betrag der Größe q abhängig. Für einen Betrag \textbar q\textbar{} \textless{} 1 konvergiert die Reihe und kann berechnet werden zu
\begin{align}
\sum\limits_{k = 0}^{\infty}q^{k} &= \frac{1}{1 - q} \label{eq:5.19}
\end{align}
Die Summe aus Gleichung \eqref{eq:5.17} ist eine geometrische Reihe, die für \textbar \(z^{-1}\)\textbar{} \textless{} 1 konvergiert und dargestellt werden kann als
\begin{align}
X(z) &= \sum\limits_{k = 0}^{\infty}\left( z^{- 1} \right)^{k} = \frac{1}{1 - z^{- 1}} = \frac{z}{z - 1} \label{eq:5.20}
\end{align}
Auch an diesem Beispiel wird eine Verschiebung um \(k_{0}\) untersucht.
\begin{align}
X(z) &= \sum\limits_{k = 0}^{\infty}{\sigma\left[ k - k_{0} \right]  z^{- k}} = \sum\limits_{k = k_{0}}^{\infty}z^{- k} = \sum\limits_{k = 0}^{\infty}z^{- (k + k_{0})} = z^{- k_{0}}  \sum\limits_{k = 0}^{\infty}z^{- k} = z^{- k_{0}}  \frac{z}{z - 1} \label{eq:5.21}
\end{align}
Die Auswertung der Beispiele Impulsfolge und Sprungfolge legt die Vermutung nahe, dass die Verschiebung einer Folge um \(k_{0}\) nach rechts generell einer Multiplikation der ursprünglichen z-Transformierten mit \(z^{- k_{0}}\ \)entspricht.

\subsubsection{Diskrete Exponentialfolge}\label{sec:diskrete-exponentialfolge}

Die diskrete Exponentialfolge ist definiert als
\begin{align}
x[ k] &= \lambda_{0}^{k}  \sigma[ k] \label{eq:5.22}
\end{align}
mit
\begin{align}
\lambda_{0} &= r_{0}  e^{j  \Omega_{0}} \label{eq:5.23}
\end{align}
Analog zur Sprungfolge ergibt sich die z-Transformierte zu
\begin{align}
X(z) &= \sum\limits_{k = 0}^{\infty}{\sigma[ k]  \lambda_{0}^{k}  z^{- k}} = \sum\limits_{k = 0}^{\infty}{\lambda_{0}^{k}  z^{- k}} = \sum\limits_{k = 0}^{\infty}\left( \frac{\lambda_{0}}{z} \right)^{k} \label{eq:5.24}
\end{align}
Diese Gleichung kann für die Berechnung der Korrespondenz für die diskrete Exponentialfolge verwendet werden. Ist der Betrag von \textbar \(λ_{0}\)\(z^{-1}\)\textbar{} \textless{} 1, konvergiert die Summe, und es gilt:
\begin{align}
X(z) &= \sum\limits_{k = 0}^{\infty}\left( \lambda_{0}  z^{- 1} \right)^{k} = \frac{1}{1 - \lambda_{0}  z^{- 1}} = \frac{z}{z - \lambda_{0}} \label{eq:5.25}
\end{align}
Die bislang berechneten z-Transformierten werden über die Definitionsgleichung berechnet. Zur Vereinfachung der Berechnung werden in Abschnitt \hyperref[_Ref372383177]{0} einige Rechenregeln der z-Transformation hergeleitet.

\subsection{Existenz der z-Transformierten}\label{sec:existenz-der-z-transformierten}

Bei der Berechnung einer z-Transformierten wird eine unendliche Reihe ausgewertet. Die\\
z-Transformierte existiert nur dann, wenn diese Reihe konvergiert. Bei der Berechnung der\\
z-Transformierten müssen teilweise Bedingungen an die Variable z gestellt werden, um die Konvergenz der Reihe sicherzustellen. Das führt zu der Frage, unter welchen Bedingungen die\\
z-Transformierte \(X(z)\) mit der Definitionsgleichung
\begin{align}
X(z) &= \sum\limits_{k = 0}^{\infty}{x[ k]  z^{- k}} \label{eq:5.26}
\end{align}
existiert. Für beidseitig begrenzte Signale \(x[k]\) mit begrenzter Amplitude konvergiert die Summe in Gleichung \eqref{eq:5.26} immer. Zur Untersuchung der Konvergenz bei nur einseitig begrenzten Folgen \(x[k]\) werden die Folgen mit Hilfe einer Exponentialfolge nach oben abgeschätzt.
\begin{align}
\left| x[ k] \right| \leq M  r^{k} \label{eq:5.27}
\end{align}
Konvergiert die resultierende geometrische Reihe, existiert auch ihre z-Transformierte \(X(z)\). Mit der Abschätzung der Folge \(x[k]\) gilt für \(X(z)\)
\begin{align}
X(z) &= \sum\limits_{k = 0}^{\infty}{x[ k]  z^{- k}} \leq \sum\limits_{k = 0}^{\infty}\left| x[ k]  z^{- k} \right| \leq \sum\limits_{k = 0}^{\infty}{M  r^{k}  \left| z^{- k} \right|} = M  \sum\limits_{k = 0}^{\infty}\left| \frac{r}{z} \right|^{k} \label{eq:5.28}
\end{align}
Damit existiert die z-Transformierte, wenn für den Betrag der Basis gilt:
\begin{align}
\left| \frac{r}{z} \right| < 1 \label{eq:5.29}
\end{align}
Diese Bedingung erfüllen alle Werte der z-Ebene, die außerhalb des Kreises mit dem Radius r liegen. Bild \ref{fig:5.1} zeigt den Konvergenzbereich in der komplexen z-Ebene.

\begin{figure}
  \centering
  \includegraphics[width=\imagewidthfactor\columnwidth, alt={Konvergenzbereich der z-Transformierten \(X(z)\) einer Folge \(x[k]\)}]{new_media/5_1_konvergenzbereich_der_ztransform.png}
  \caption{Konvergenzbereich der z-Transformierten \(X(z)\) einer Folge \(x[k]\)}
  \label{fig:5.1}
\end{figure}

In praktischen Anwendungen können unendliche Signalfolgen mit einer Exponentialfolge nach oben abgeschätzt werden, sodass die Konvergenz in den meisten Fällen gegeben ist.

\subsection{Zusammenhang zwischen z-Transformation und Laplace-Transformation}\label{sec:5.1.4}

Mit Hilfe der Laplace-Transformation lassen sich einige Signaleigenschaften bestimmen. Um an diese Interpretationsmöglichkeiten anzuknüpfen, wird ein Zusammenhang zwischen der s-Ebene der Laplace-Transformation und der z-Ebene der z-Transformation hergestellt. Bei der Herleitung der z-Transformation wird die Substitution
\begin{align}
z &= e^{T_{A}  s} \label{eq:5.30}
\end{align}
verwendet. Mit dieser Substitution wird die komplexe s-Ebene in die komplexe z-Ebene abgebildet. Bild \ref{fig:5.2} zeigt die beiden Ebenen.

\begin{figure}
  \centering
  \includegraphics[width=\imagemaxwidthfactor\columnwidth, alt={Darstellung von Bereichen in der s-Ebene und deren transformierten Flächen in der z-Ebene}]{new_media/5_2_darstellung_von_bereichen_in_der.png}
  \caption{Darstellung von Bereichen in der s-Ebene und deren transformierten Flächen in der z-Ebene}
  \label{fig:5.2}
\end{figure}

Die imaginäre Achse der s-Ebene s = jω wird abgebildet in die Variable
\begin{align}
z &= e^{j  \omega  T_{A}} \label{eq:5.31}
\end{align}
Das entspricht dem Einheitskreis, der für - ∞ \textless{} ω \textless{} ∞ periodisch in 2π durchlaufen wird. Die imaginäre Achse der s-Ebene wird demnach auf den Einheitskreis der z-Ebene abgebildet. Bei der Abtastung von Signalen entsteht ein Spektrum, das periodisch in \(ω_{A}\) ist. Das Basisband reicht dabei von - \(ω_{A}\)/2 bis +~\(ω_{A}\)/2, so dass der Einheitskreis genau einmal durchlaufen wird. Alle periodischen Wiederholungen des Spektrums werden mit der z-Transformation auf dem Einheitskreis übereinander abgebildet. Deshalb eignet sich die z-Transformation besonders für die Beschreibung abgetasteter Signale.

Eine Linie der s-Ebene mit konstantem Realteil \(δ_{0}\) wird auf einen Kreis abgebildet, der einen Radius
\begin{align}
|z| &= \left| e^{\left( \delta_{0} + j  \omega \right)  T_{A}} \right| = \left| e^{\delta_{0}  T_{A}} \right|  \left| e^{j  \omega  T_{A}} \right| = e^{\delta_{0}  T_{A}} \label{eq:5.32}
\end{align}
aufweist. Mit negativem Realteil ergibt sich eine Lage innerhalb des Einheitskreises, mit positivem Realteil eine Lage um den Einheitskreis herum. Daraus folgt, dass die gesamte linke s-Halbebene in den Einheitskreis abgebildet wird, während die rechte s-Halbebene um den Einheitskreis herum abgebildet wird. Tabelle \ref{tab:5.1} fasst korrespondierende Elemente der s- und z-Ebene zusammen.

\begin{application}{Transformation}
  Im Online-Portal \emph{Systemtheorie Online} verdeutlicht die Applikation \emph{Transformation} den Zusammenhang zwischen Punkten in der Laplace-Ebene und Punkten in der z-Ebene.
  \tcblower
  \url{h-ka.de/eit/digitale-signalverarbeitung/übungen/kapitel2}
\end{application}


\begin{summary}{Tabellarische Übersicht über korrespondierende Elemente der s- und z-Ebene}
\label{tab:5.1}
\begin{center}
  \begin{tabular}{ll}
    \toprule
    s-Ebene & z-Ebene \\
    \midrule
    imaginäre Achse & Einheitskreis \\
    linke komplexe Ebene & Inneres des Einheitskreises \\
    rechte komplexe Ebene & Äußeres des Einheitskreises \\
    Ursprung \(s = 0\) & \(z = 1\) \\
    halbe Abtastfrequenz jω = jω\textsubscript{A}/2 ± jkω\textsubscript{A} & z = - 1 \\
    \bottomrule
  \end{tabular}
\end{center}
\end{summary}

Der Zusammenhang zwischen s- und z-Ebene wird bei der Diskussion der Signal- und Systemeigenschaften in den folgenden Abschnitten wieder aufgegriffen.

\subsection{Pollage und komplexe Exponential-Folge}\label{sec:pollage-und-komplexe-exponential-folge}

Im vorangegangenen Abschnitt wird die z-Transformierte der Exponentialfolge
\begin{align}
x[ k] &= \lambda_{0}^{k}  \sigma[ k] \label{eq:5.33}
\end{align}
berechnet zu

\begin{table}
  \centering
  \caption{Zusammenhang zwischen Pollage und Signalfolge bei reellen Polen}
  \label{tab:5.2}
  \begin{tabular}{ll}
    \toprule
    \noalign{}

\[X(z) = \frac{z}{z - \lambda_{0}}\]
 & 
\protect\phantomsection\label{_Ref273456945}{}(5.34) \\
    \bottomrule
  \end{tabular}
\end{table}

{\def\LTcaptype{none} % do not increment counter
\begin{table}
  \centering
  \caption{Tabellarische Übersicht}
  \begin{tabular}{ll}
    \toprule
    Pollage \(X(z)\) & Signalfolge \(x[k]\) \\
    \midrule
    \includegraphics[width=\imagewidthfactor\columnwidth, alt={Ein Bild, das Reihe, Kreis, Screenshot, Diagramm enthält. Automatisch generierte Beschreibung}]{new_media/5_3_tabellarische_bersicht.png} & \includegraphics[width=\imagewidthfactor\columnwidth, alt={Ein Bild, das Screenshot, Reihe enthält. Automatisch generierte Beschreibung}]{new_media/5_4_tabellarische_bersicht.png} \\
    \includegraphics[width=\imagewidthfactor\columnwidth, alt={Ein Bild, das Reihe, Kreis, Screenshot, Diagramm enthält. Automatisch generierte Beschreibung}]{new_media/5_5_tabellarische_bersicht.png} & \includegraphics[width=\imagewidthfactor\columnwidth, alt={Ein Bild, das Screenshot, Reihe enthält. Automatisch generierte Beschreibung}]{new_media/5_6_tabellarische_bersicht.png} \\
    \includegraphics[width=\imagewidthfactor\columnwidth, alt={Ein Bild, das Reihe, Screenshot, Diagramm, Kreis enthält. Automatisch generierte Beschreibung}]{new_media/5_7_tabellarische_bersicht.png} & \includegraphics[width=\imagewidthfactor\columnwidth, alt={Ein Bild, das Screenshot, Reihe enthält. Automatisch generierte Beschreibung}]{new_media/5_8_tabellarische_bersicht.png} \\
    \includegraphics[width=\imagewidthfactor\columnwidth, alt={Ein Bild, das Screenshot, Reihe, Kreis, Diagramm enthält. Automatisch generierte Beschreibung}]{new_media/5_9_tabellarische_bersicht.png} & \includegraphics[width=\imagewidthfactor\columnwidth, alt={Ein Bild, das Reihe, Screenshot, Diagramm, Rechteck enthält. Automatisch generierte Beschreibung}]{new_media/5_10_tabellarische_bersicht.png} \\
    \includegraphics[width=\imagewidthfactor\columnwidth, alt={Ein Bild, das Reihe, Kreis, Screenshot, Diagramm enthält. Automatisch generierte Beschreibung}]{new_media/5_11_tabellarische_bersicht.png} & \includegraphics[width=\imagewidthfactor\columnwidth, alt={Ein Bild, das Reihe, Diagramm, Rechteck, Screenshot enthält. Automatisch generierte Beschreibung}]{new_media/5_12_tabellarische_bersicht.png} \\
    \includegraphics[width=\imagewidthfactor\columnwidth, alt={Ein Bild, das Reihe, Kreis, Diagramm, Screenshot enthält. Automatisch generierte Beschreibung}]{new_media/5_13_tabellarische_bersicht.png} & \includegraphics[width=\imagewidthfactor\columnwidth, alt={Ein Bild, das Reihe, Screenshot, Diagramm, Rechteck enthält. Automatisch generierte Beschreibung}]{new_media/5_14_tabellarische_bersicht.png} \\
    \bottomrule
  \end{tabular}
\end{table}
}

Aus Gleichung \eqref{eq:5.34} kann der zu der komplexen Exponentialfolge zugehörige Pol \(λ_{0}\) in der komplexen z-Ebene abgelesen werden. Liegt der Pol auf der reellen Achse, handelt es sich um eine reelle Folge, die bei negativem Realteil alternierend ist. Für einen Betrag von \textbar \(λ_{0}\)\textbar{} \textless{} 1 ist die Folge konvergent, für\\
\textbar \(λ_{0}\)\textbar{} \textgreater{} 1 ist die Folge divergent und für \textbar \(λ_{0}\)\textbar{} = 1 bleibt der Betrag der Folge konstant. Die Lage eines Poles in der z-Ebene kann damit einem Signalverhalten zugeordnet werden, das in Tabelle \ref{tab:5.2} skizziert wird.

Liegen z-Transformierte mit konjugiert komplexem Polpaar der Form
\begin{align}
X(z) &= \frac{z}{z - r_{0}  e^{j  \Omega_{0}}} + \frac{z}{z - r_{0}  e^{- j  \Omega_{0}}} \label{eq:5.35}
\end{align}
vor, handelt es sich im Zeitbereich um Folgen der Form
\begin{align}
x[ k] &= r_{0}^{k}  \left( e^{j  \Omega_{0}  k} + e^{- j  \Omega_{0}  k} \right)  \sigma[ k] = 2  r_{0}^{k}  \cos\left( \Omega_{0}  k \right)  \sigma[ k] \label{eq:5.36}
\end{align}
Für einen Betrag von \textbar \(r_{0}\)\textbar{} \textless{} 1 ist die Folge konvergent, für \textbar \(r_{0}\)\textbar{} \textgreater{} 1 ist die Folge divergent und für \textbar \(r_{0}\)\textbar{} = 1 bleibt der Betrag der Folge konstant. Die Frequenz steigt mit wachsender Phase \(Ω_{0}\) des Pols, das Maximum wird erreicht für Pole mit einer Phase \(Ω_{0}\) = ± π. Diese Interpretation entspricht bei entsprechender Umrechnung der Pole dem Zusammenhang zwischen der Pollage in der z-Ebene und der zugehörigen Zeitfunktion. Die Lage eines Poles in der z-Ebene kann damit einem Signalverhalten zugeordnet werden, das in Tabelle \ref{tab:5.3} skizziert wird.

\begin{table}
  \centering
  \caption{Zusammenhang zwischen Pollage und Signalfolge bei konjugiert komplexen Polpaaren}
  \label{tab:5.3}
  \begin{tabular}{ll}
    \toprule
    Pollage \(X(z)\) & Signalfolge \(x[k]\) \\
    \midrule
    \includegraphics[width=\imagewidthfactor\columnwidth, alt={Ein Bild, das Reihe, Screenshot, Diagramm, Kreis enthält. Automatisch generierte Beschreibung}]{new_media/5_15_zusammenhang_zwischen_pollage_u.png} & \includegraphics[width=\imagewidthfactor\columnwidth, alt={Ein Bild, das Screenshot, Reihe, Rechteck, Diagramm enthält. Automatisch generierte Beschreibung}]{new_media/5_16_zusammenhang_zwischen_pollage_u.png} \\
    \includegraphics[width=\imagewidthfactor\columnwidth, alt={Ein Bild, das Reihe, Diagramm, Kreis, Screenshot enthält. Automatisch generierte Beschreibung}]{new_media/5_17_zusammenhang_zwischen_pollage_u.png} & \includegraphics[width=\imagewidthfactor\columnwidth, alt={Ein Bild, das Reihe, Diagramm, Screenshot, Rechteck enthält. Automatisch generierte Beschreibung}]{new_media/5_18_zusammenhang_zwischen_pollage_u.png} \\
    \includegraphics[width=\imagewidthfactor\columnwidth, alt={Ein Bild, das Reihe, Diagramm, Screenshot, Kreis enthält. Automatisch generierte Beschreibung}]{new_media/5_19_zusammenhang_zwischen_pollage_u.png} & \includegraphics[width=\imagewidthfactor\columnwidth, alt={Ein Bild, das Reihe, Screenshot, Diagramm, Rechteck enthält. Automatisch generierte Beschreibung}]{new_media/5_20_zusammenhang_zwischen_pollage_u.png} \\
    \bottomrule
  \end{tabular}
\end{table}

\begin{application}{Komplexe Exponentialfolge}
  Im Online-Portal \emph{Systemtheorie Online} verdeutlicht die Applikation \emph{Komplexe Exponentialfolge} den Zusammenhang zwischen der Pollage in der komplexen Ebene und dem Verhalten der Signalfolge.
  \tcblower
  \url{h-ka.de/eit/digitale-signalverarbeitung/übungen/kapitel2}
\end{application}


\protect\phantomsection\label{_Ref372383177}{}

\section{Rechenregeln der z-Transformation}\label{sec:5.2}

Wie bei den bisher vorgestellten Transformationen besitzt die z-Transformation Rechenregeln, die das Rechnen und die Berechnung weiterer Korrespondenzen vereinfachen. Sie werden hergeleitet und ihr Nutzen an Beispielen verdeutlicht.

\subsection{Linearität}\label{sec:linearituxe4t-2}

Wie die Laplace- und die Fourier-Transformation ist auch die z-Transformation eine lineare Transformation. Durch Einsetzen in die Definitionsgleichung
\begin{align}
Z\left\{ \nu_{1}  x_{1}[ k] + \nu_{2}  x_{2}[ k] \right\} &= \sum\limits_{k = 0}^{\infty}{\left( \nu_{1}  x_{1}[ k] + \nu_{2}  x_{2}[ k] \right)  z^{- k}} \label{eq:5.37}
\end{align}
und durch Anwendung des Distributivgesetzes
\begin{align}
\sum\limits_{k = 0}^{\infty}{\left( \nu_{1}  x_{1}[ k] + \nu_{2}  x_{2}[ k] \right)  z^{- k}} &= \sum\limits_{k = 0}^{\infty}{\nu_{1}  x_{1}[ k]  z^{- k}} + \sum\limits_{k = 0}^{\infty}{\nu_{2}  x_{2}[ k]  z^{- k}} = \nu_{1}  X_{1}(z) + \nu_{2}  X_{2}(z) \label{eq:5.38}
\end{align}
wird die Linearitätseigenschaft bewiesen.

\begin{example}{Linearität}
\label{example:beispiel-linearituxe4t}

Gesucht wird die z-Transformierte der Folge \(x[k]\)
\begin{align*}
x[ k] &= 2  \sigma[ k] + 5  3^{k}  \sigma[ k] 
\end{align*}
Die z-Transformierten der Sprung- und Exponentialfolge sind bekannt. Damit ergibt sich \(X(z)\) zu
\begin{align*}
X(z) &= 2  \frac{z}{z - 1} + 5  \frac{z}{z - 3} 
\end{align*}
\end{example}

\subsection{Verschiebung der Folge nach rechts}\label{sec:verschiebung-der-folge-nach-rechts}

Bereits bei der Herleitung der verschobenen Impuls- und Sprungfolge wird die Verschiebungsregel der z-Transformation angedeutet. Die allgemeine Herleitung für eine Zeitverschiebung nach rechts ergibt sich für den Fall, dass \(x[k]\) eine kausale Folge ist, aus
\begin{align}
{\boldsymbol{\mathit{Z}}\left\{x\left[k-k_{0} \right]\right\}} &=  {\sum\limits _{k=0}^{\infty }x\left[k-k_{0} \right] z^{-k}  =\sum\limits _{k=k_{0} }^{\infty }x\left[k-k_{0} \right] z^{-k}  =\sum\limits _{k=0}^{\infty }x[k] z^{-(k+k_{0} )}  =z^{-k_{0} }  \sum\limits _{k=0}^{\infty }x[k] z^{-k}  } \\
{} &=  {z^{-k_{0} }  X(z)}
\label{eq:5.41}
\end{align}
Diskrete Systeme lassen sich über Differenzengleichungen beschreiben. Aus diesem Grund ist die Verschiebungsregel von vergleichbarer Bedeutung wie die Differentiationsregel der Laplace-Transformation.

\begin{example}{Verschiebung nach rechts}
\label{example:beispiel-verschiebung-nach-rechts}

Eine Rechteckfolge von 0 bis 4 kann im Zeitbereich dargestellt werden als
\begin{align*}
x[ k] &= \sigma[ k] - \sigma[ k - 5] 
\end{align*}
Daraus ergibt sich mit der Verschiebungsregel die z-Transformierte zu
\begin{align*}
X(z) &= \frac{z}{z - 1}  \left( 1 - z^{- 5} \right) 
\end{align*}
\end{example}

\subsection{Verschiebung der Folge nach links}\label{sec:verschiebung-der-folge-nach-links}

Die Regel zur Verschiebung nach links kann ähnlich hergeleitet werden wie die Regel zur Verschiebung nach rechts.
\begin{align}
Z\left\{ x[ k + k_{0}] \right\} &= \sum\limits_{k = 0}^{\infty}{x[ k + k_{0}]  z^{- k}} = z^{k_{0}}  \sum\limits_{k = 0}^{\infty}{x[ k + k_{0}]  z^{- (k + k_{0})}} = z^{k_{0}}  \sum\limits_{k = k_{0}}^{\infty}{x[ k]  z^{- k}} \label{eq:5.44}
\end{align}
Um die Summe bei \(k = 0\) beginnen zu lassen, müssen einige Summanden addiert und wieder subtrahiert werden. Damit ergibt sich die Verschiebungsregel für eine Verschiebung nach links.
\begin{align}
{\boldsymbol{\mathit{Z}}\left\{x\left[k+k_{0} \right]\right\}} &=  {z^{k_{0} }  \sum\limits _{k=0}^{k_{0} -1}x\left[k\right] z^{-k}  +z^{k_{0} }  \sum\limits _{k=k_{0} }^{\infty }x\left[k\right] z^{-k}  -z^{k_{0} }  \sum\limits _{k=0}^{k_{0} -1}x\left[k\right] z^{-k}  } \\
{} &=  {z^{k_{0} }  \sum\limits _{k=0}^{\infty }x\left[k\right] z^{-k}  -z^{k_{0} }  \sum\limits _{k=0}^{k_{0} -1}x\left[k\right] z^{-k}  =z^{k_{0} }  X\left(z\right)-z^{k_{0} }  \sum\limits _{k=0}^{k_{0} -1}x\left[k\right] z^{-k}  }
\label{eq:5.45}
\end{align}
Bei der Diskussion von Systemen im z-Bereich wird sich zeigen, dass sich diese Verschiebungsregel zur Lösung von Differenzengleichungen mit Anfangsbedingungen eignet.

Beispiel: Verschiebung nach links

Die Differenzengleichung
\begin{align}
x[ k] - 0.5  x[ k - 1] &= 0 \label{eq:5.46}
\end{align}
soll für den Anfangswert \(x[0]\) = 5 gelöst werden. Dazu wird die Differenzengleichung umgeformt zu
\begin{align}
x[ k + 1] - 0.5  x[ k] &= 0 \label{eq:5.47}
\end{align}
und die Verschiebungsregel angewendet.
\begin{align}
z^{1}  X(z) - x[ 0]  z^{1} + 0.5  X(z) &= 0 \label{eq:5.48}
\end{align}
Auflösen nach \(X(z)\) führt zu der z-Transformierten
\begin{align}
X(z) &= \frac{z  x[ 0]}{z + 0.5} = \frac{z  5}{z + 0.5} \label{eq:5.49}
\end{align}
Mit Hilfe der bereits diskutierten Rechenregeln kann die z-Transformierte in den Zeitbereich zurücktransformiert werden.
\begin{align}
x[ k] &= 5  ( - 0.5)^{k}  \sigma[ k] \label{eq:5.50}
\end{align}
▪

\subsection{Modulation}\label{sec:modulation}

Eine Modulation einer Folge \(x[k]\) mit der Exponentialfolge \(λ^{k}\) lässt sich im z-Bereich darstellen als
\begin{align}
Z\left\{ \lambda^{k}  x[ k] \right\} &= \sum\limits_{k = 0}^{\infty}{\lambda^{k}  x[ k]  z^{- k}} = \sum\limits_{k = 0}^{\infty}{x[ k]  \left( \frac{z}{\lambda} \right)^{- k}} = X\left( \frac{z}{\lambda} \right) \label{eq:5.51}
\end{align}
\begin{example}{Modulation}
\label{example:beispiel-modulation}

Die Modulationsregel kann zum Nachweis der z-Transformierten der Exponentialfolge angewendet werden. Die z-Transformierte der Sprungfolge ist bekannt. Die Exponentialfolge ergibt sich aus
\begin{align*}
x[ k] &= \lambda^{k}  \sigma[ k] 
\end{align*}
Nach der Modulationsregel ist damit die z-Transformierte der Exponentialfolge
\begin{align*}
X(z) &= \left. \ \frac{z}{z - 1} \right|_{z = \frac{z}{\lambda}} = \frac{\frac{z}{\lambda}}{\frac{z}{\lambda} - 1} = \frac{z}{z - \lambda} 
\end{align*}
\end{example}

\subsection{Lineare Gewichtung}\label{sec:lineare-gewichtung}

Die Eigenschaft der linearen Gewichtung ergibt sich durch Ableitung der z-Transformierten.
\begin{align}
\frac{dX}{dz} &= \frac{d}{dz}\sum\limits_{k = 0}^{\infty}{x[ k]  z^{- k}} = - k  \sum\limits_{k = 0}^{\infty}{x[ k]  z^{- k - 1}} = - z^{- 1}  \sum\limits_{k = 0}^{\infty}{k  x[ k]  z^{- k}} \label{eq:5.54}
\end{align}
Multiplikation der Gleichung mit - z führt zu der gesuchten Rechenregel der linearen Gewichtung.
\begin{align}
Z\left\{ k  x[ k] \right\} &= \sum\limits_{k = 0}^{\infty}{k  x[ k]  z^{- k}} = - z  \frac{dX}{dz} \label{eq:5.55}
\end{align}
\begin{example}{Gewichtung}
\label{example:beispiel-gewichtung}

Die lineare Gewichtung wird dazu verwendet, die z-Transformierte der Rampenfolge
\begin{align*}
x[ k] &= k  \sigma[ k] 
\end{align*}
zu berechnen. Mit der Regel für die lineare Gewichtung und der z-Transformierten der Sprungfolge ergibt sich
\begin{align*}
X(z) &= - z  \frac{d}{dz}\frac{z}{z - 1} = - z  \frac{z - 1 - z}{(z - 1)^{2}} = \frac{z}{(z - 1)^{2}} 
\end{align*}
\end{example}

\subsection{Differenz}\label{sec:differenz}

Die Differenz zweier aufeinanderfolgender Werte entspricht der Ableitung im zeitkontinuierlichen Bereich. Ihre z-Transformierte ergibt sich direkt aus dem Verschiebungssatz zu
\begin{align}
Z\left\{ x[ k] - x[ k - 1] \right\} &= \left( 1 - z^{- 1} \right)  X(z) = \frac{z - 1}{z}  X(z) \label{eq:5.58}
\end{align}
\begin{example}{Differenz}
\label{example:beispiel-differenz}

Die Impulsfolge kann dargestellt werden als Differenz zweier Sprungfolgen:
\begin{align*}
\delta[ k] &= \sigma[ k] - \sigma[ k - 1] 
\end{align*}
Damit ergibt sich ihre z-Transformierte zu
\begin{align*}
X(z) &= \frac{z - 1}{z}  \frac{z}{z - 1} = 1 
\end{align*}
Das Ergebnis stimmt mit der Berechnung über die Definitionsgleichung überein.

\end{example}

\subsection{Summation}\label{sec:summation}

Dem Integral für zeitkontinuierliche Funktionen entspricht bei Folgen die Summation. Zum Beweis der Summationsregel wird die Folge \(x[k]\) dargestellt als Differenz zweier Summen:
\begin{align}
x[ k] &= \sum\limits_{\kappa = 0}^{k}{x[\kappa] -}\sum\limits_{\kappa = 0}^{k - 1}{x[\kappa]} \label{eq:5.61}
\end{align}
Die z-Transformation der Gleichung ergibt
\begin{align}
X(z) &= \left( 1 - z^{- 1} \right)  Z\left( \sum\limits_{\kappa = 0}^{k}{x[\kappa]} \right) = \frac{z - 1}{z}  Z\left( \sum\limits_{\kappa = 0}^{k}{x[\kappa]} \right) \label{eq:5.62}
\end{align}
Durch Umstellen ergibt sich die Summationsregel der z-Transformation
\begin{align}
Z\left\{ \sum\limits_{\kappa = 0}^{k}{x[\kappa]} \right\} &= \frac{z}{z - 1}  X(z) \label{eq:5.63}
\end{align}
\begin{example}{Summation}
\label{example:beispiel-summation}

Die Sprungfolge kann als Summe über die Impulsfolge dargestellt werden.
\begin{align*}
\sigma[ k] &= \sum\limits_{\kappa = 0}^{k}{\delta[\kappa]} 
\end{align*}
Entsprechend ist der Zusammenhang zwischen den beiden z-Transformierten
\begin{align*}
Z\left\{ \sigma[ k] \right\} &= Z\left\{ \sum\limits_{\kappa = 0}^{k}{\delta[\kappa]} \right\} = 1  \frac{z}{z - 1} = \frac{z}{z - 1} 
\end{align*}
Auch die Rampenfolge lässt sich als Summe über die Sprungfolge beschreiben. Dabei ist zu beachten, dass die Rampenfolge für \(k = 0\) den Wert 0 besitzt. Damit ergibt sich für die Summe der Startwert\\
\(k = 1\) und die Summenformel der Rampenfolge lautet
\begin{align*}
x[ k] &= \sum\limits_{\kappa = 0}^{k}{\sigma[\kappa - 1]} 
\end{align*}
Mit der Summations- und Verschiebungsregel ergibt sich die z-Transformierte zu
\begin{align*}
X(z) &= \frac{z}{z - 1}  z^{- 1}  \frac{z}{z - 1} = \frac{z}{(z - 1)^{2}} 
\end{align*}
Das Ergebnis stimmt mit der Berechnung über die lineare Gewichtung überein.

\end{example}

\subsection{Multiplikation}\label{sec:multiplikation}

Die Multiplikationsregel wird im folgenden Abschnitt mit Hilfe der Rücktransformation hergeleitet. Sie wird hier der Vollständigkeit halber aufgeführt.
\begin{align}
Z\left\{ x[ k]  y[ k] \right\} &= \frac{1}{2  \pi  j}  \oint_{}^{}{Y\left( \frac{z}{v} \right)  X(v)  v^{- 1}\ dv} \label{eq:5.68}
\end{align}
\subsection{Faltung}\label{sec:faltung}

Der Faltungssatz ist für die Berechnung von Systemantworten von großer Bedeutung und wird ausführlich in Kapitel 8 behandelt. Die z-Transformierte der Faltung zweier Folgen ergibt sich durch Einsetzen in die Definitionsgleichung der z-Transformation zu
\begin{align}
Z\left\{ g[ k]*x[ k] \right\} &= Z\left\{ \sum\limits_{\kappa = 0}^{\infty}{g[\kappa]  x[ k - \kappa]} \right\} = \sum\limits_{k = 0}^{\infty}{\sum\limits_{\kappa = 0}^{\infty}{g[\kappa]  x[ k - \kappa]}}  z^{- k} \label{eq:5.69}
\end{align}
Vertauschen der Summationsreihenfolge und Substitution ergibt
\begin{align}
{\sum\limits _{k=0}^{\infty }\sum\limits _{\kappa =0}^{\infty }x\left[\kappa \right] g\left[k-\kappa \right]   z^{-k} } &=  {\sum\limits _{\kappa =0}^{\infty }\sum\limits _{k=0}^{\infty }x\left[\kappa \right] g\left[k-\kappa \right]   z^{-k} } \\
{} &=  {\sum\limits _{\kappa =0}^{\infty }x\left[\kappa \right] \sum\limits _{k=0}^{\infty }g\left[k-\kappa \right]   z^{-k}  z^{\kappa }  z^{-\kappa } } \\
{} &=  {\sum\limits _{\kappa =0}^{\infty }x\left[\kappa \right] z^{-\kappa }  \sum\limits _{k=0}^{\infty }g\left[k-\kappa \right]   z^{-\left(k-\kappa \right)} } \\
{} &=  {\sum\limits _{\kappa =0}^{\infty }x\left[\kappa \right] z^{-\kappa }  \sum\limits _{n=0}^{\infty }g \left[n\right] z^{-n}   =G\left(z\right) X\left(z\right)}
\label{eq:5.70}
\end{align}
Diese Rechenregel wird bei der Diskussion zeitdiskreter Systeme im Bildbereich aufgegriffen.

\subsection{Anfangswertsatz}\label{sec:anfangswertsatz}

Analog zu den Grenzwertsätzen der Laplace-Transformation ergeben sich Anfangs- und Endwertsätze der z-Transformation. Die Definitionsgleichung der z-Transformation lautet
\begin{align}
X(z) &= \sum\limits_{k = 0}^{\infty}{x[ k]  z^{- k}} = x[ 0] + x[ 1]  z^{- 1} + x[ 2]  z^{- 2} + ... \label{eq:5.71}
\end{align}
Im Grenzfall z → ∞ werden alle Summanden, die einen Faktor \(z^{-k}\) mit k \textgreater{} 0 besitzen, zu null. Übrig bleibt demnach der Ausdruck
\begin{align}
x[ 0] &= \lim_{z \rightarrow \infty}X(z) \label{eq:5.72}
\end{align}
Da bei der Herleitung des Anfangswertsatzes keine Einschränkung gemacht wird, kann diese Regel immer angewendet werden.

\begin{example}{Anfangswertsatz}
\label{example:beispiel-anfangswertsatz}

Die kausale Exponentialfolge besitzt die z-Transformierte
\begin{align*}
X(z) &= \frac{z}{z - \lambda_{0}} 
\end{align*}
Für den Grenzübergang z → ∞ ergibt sich der Anfangswert zu
\begin{align*}
\lim_{z \rightarrow \infty}X(z) &= \lim_{z \rightarrow \infty}\frac{z}{z - \lambda_{0}} = \lim_{z \rightarrow \infty}\frac{1}{1 - \frac{\lambda_{0}}{z}} = 1 
\end{align*}
Das Ergebnis stimmt mit dem Wert \(x[0]\) im Zeitbereich überein.

\end{example}

\subsection{Endwertsatz}\label{sec:endwertsatz}

Zur Berechnung des Grenzwertes der Folge \(x[k]\) für k → ∞ wird vorausgesetzt, dass dieser Grenzwert existiert. Zur Herleitung des Endwertsatzes wird die z-Transformierte \(X(z)\) des Signals \(x[k]\) dargestellt als ein Produkt zweier z-Transformierter
\begin{align}
X(z) &= \frac{z}{z - 1}  Y(z) \label{eq:5.75}
\end{align}
Diese Gleichung kann mit der Faltungsregel in den Zeitbereich transformiert werden.
\begin{align}
x[ k] &= \sum\limits_{\kappa = 0}^{\infty}{y[ n]  \sigma[ k - \kappa]} = \sum\limits_{\kappa = 0}^{k}{y[\kappa]} \label{eq:5.76}
\end{align}
Für den Grenzwert k → ∞ kann die Summe dargestellt werden als
\begin{align}
\lim_{k \rightarrow \infty}x[ k] &= \sum\limits_{\kappa = 0}^{\infty}{y[\kappa]} = \sum\limits_{\kappa = 0}^{\infty}{y[\kappa]}  1^{- \kappa} = Y(1) = \lim_{z \rightarrow 1}\frac{z - 1}{z}  X(z) = \lim_{z \rightarrow 1}(z - 1)  X(z) \label{eq:5.77}
\end{align}
Mit der Definitionsgleichung der z-Transformation
\begin{align}
Y(z) &= \sum\limits_{k = 0}^{\infty}{y[ k]  z^{- k}} \label{eq:5.78}
\end{align}
ergibt sich für \(z = 1\) der Endwertsatz der z-Transformation
\begin{align}
\lim_{k \rightarrow \infty}x[ k] &= \sum\limits_{\kappa = 0}^{\infty}{h[\kappa]}  1^{- \kappa} = H(1) = \lim_{z \rightarrow 1}\frac{z - 1}{z}  X(z) = \lim_{z \rightarrow 1}(z - 1)  X(z) \label{eq:5.79}
\end{align}
\begin{example}{Endwertsatz}
\label{example:beispiel-endwertsatz}

Die kausale Exponentialfolge besitzt die z-Transformierte
\begin{align*}
X(z) &= \frac{z}{z - \lambda_{0}} 
\end{align*}
Mit dem Grenzübergang z → 1 des Ausdrucks
\begin{align*}
\lim_{z \rightarrow 1}(z - 1)  \frac{z}{z - \lambda_{0}} &= 0 
\end{align*}
ergibt sich der Endwert von \(x[k]\) für k → ∞ zu null. Das Ergebnis stimmt mit dem Grenzwert im Zeitbereich überein, wenn die Basis \(λ_{0}\) einen Betrag \textbar \(λ_{0}\)\textbar{} \textless{} 1 aufweist, also wenn ein Grenzwert existiert.

\end{example}

\subsection{Zusammenfassung der Rechenregeln zur z-Transformation}\label{sec:zusammenfassung-der-rechenregeln-zur-z-transformation}

Zur besseren Übersicht stellt Tabelle \ref{tab:5.4} die wesentlichen Eigenschaften der z-Transformation noch einmal zusammen.

\begin{table}
  \centering
  \caption{Rechenregeln der z-Transformation}
  \label{tab:5.4}
  \begin{tabular}{lll}
    \toprule
    Regel & Folge \(x[k]\) & z-Transformierte \(X(z)\) \\
    \midrule
    Linearität & \(\nu_{1}  x_{1}\lbrack k\rbrack + \nu_{2}  x_{2}\lbrack k\rbrack\) & \(\nu_{2}  X_{1}(z) + \nu_{2}  X_{2}(z)\) \\
    Zeitverschiebung nach rechts & \(x\left\lbrack k - k_{0} \right\rbrack\ für\ k_{0} > 0\) & \(z^{- k_{0}}  X(z)\) \\
    Zeitverschiebung nach links & \(x\lbrack k + k_{0}\rbrack\ für\ k_{0} > 0\) & \(z^{k_{0}}  X(z) - z^{k_{0}}  \sum\limits_{n = 0}^{k_{0} - 1}{x\lbrack n\rbrack  z^{- n}}\) \\
    Modulation & \(\lambda^{k}  x\lbrack k\rbrack\) & \(X\left( \frac{z}{\lambda} \right)\) \\
    lineare Gewichtung & \(k  x\lbrack k\rbrack\) & \(- z  \frac{d}{dz}  X(z)\) \\
    Differenz & \(x\lbrack k\rbrack - x\lbrack k - 1\rbrack\) & \(\frac{z - 1}{z}  X(z)\) \\
    Summation & \(\sum\limits_{\kappa = 0}^{k}{x\lbrack\kappa\rbrack}\) & \(\frac{z}{z - 1}  X(z)\) \\
    Multiplikation & \(x_{1}\lbrack k\rbrack  x_{2}\lbrack k\rbrack\) & \(\frac{1}{2  \pi  j}  \oint_{}^{}{X_{1}\left( \frac{z}{v} \right)  X_{2}(v)  v^{- 1}\ dv}\) \\
    Faltung & \(g\lbrack k\rbrack*x\lbrack k\rbrack\) & \(G(z)  X(z)\) \\
    Anfangswert & \(x\lbrack 0\rbrack\) & \(\lim_{z \rightarrow \infty}X(z)\) \\
    Endwert & \(\lim_{k \rightarrow \infty}x\lbrack k\rbrack\) & \(\lim_{z \rightarrow 1}\left( (z - 1)  X(z) \right)\) \\
    \bottomrule
  \end{tabular}
\end{table}

\subsection{Korrespondenzen der z-Transformation}\label{sec:korrespondenzen-der-z-transformation}

Die Rechenregeln zur z-Transformation erlauben die Berechnung weiterer Korrespondenzen. Tabelle \ref{tab:5.5} stellt wichtige Korrespondenzen der z-Transformation zusammen.

\begin{table}
  \centering
  \caption{Korrespondenzen der z-Transformation 1/2}
  \label{tab:5.5}
  \begin{tabular}{llllll}
    \toprule
    Nr. & Folge \(x[k]\) & \multicolumn{2}{l}{%
z-Transformierte \(X(z)\)} & \multicolumn{2}{>{\raggedright\arraybackslash}p{(\linewidth - 10\tabcolsep) * \real{0.3112} + 2\tabcolsep}@{}}{%
Konvergenzbereich} \\
    \midrule
    1 & \(\delta\lbrack k\rbrack\) & \multicolumn{2}{l}{%
\(1\)} & \multicolumn{2}{>{\raggedright\arraybackslash}p{(\linewidth - 10\tabcolsep) * \real{0.3112} + 2\tabcolsep}@{}}{%
\(beliebig\)} \\
    2 & \(\sigma\lbrack k\rbrack\) & \multicolumn{2}{l}{%
\(\frac{z}{z - 1}\)} & \multicolumn{2}{>{\raggedright\arraybackslash}p{(\linewidth - 10\tabcolsep) * \real{0.3112} + 2\tabcolsep}@{}}{%
\(|z| > 1\)} \\
    3 & \(k  \sigma\lbrack k\rbrack\) & \multicolumn{2}{l}{%
\(\frac{z}{(z - 1)^{2}}\)} & \multicolumn{2}{>{\raggedright\arraybackslash}p{(\linewidth - 10\tabcolsep) * \real{0.3112} + 2\tabcolsep}@{}}{%
\(|z| > 1\)} \\
    4 & \(k^{2}  \sigma\lbrack k\rbrack\) & \multicolumn{2}{l}{%
\(\frac{z  (z + 1)}{(z - 1)^{3}}\)} & \multicolumn{2}{>{\raggedright\arraybackslash}p{(\linewidth - 10\tabcolsep) * \real{0.3112} + 2\tabcolsep}@{}}{%
\(|z| > 1\)} \\
    5 & \(\lambda^{k}  \sigma\lbrack k\rbrack\) & \multicolumn{2}{l}{%
\(\frac{z}{z - \lambda}\)} & \multicolumn{2}{>{\raggedright\arraybackslash}p{(\linewidth - 10\tabcolsep) * \real{0.3112} + 2\tabcolsep}@{}}{%
\({|z| > |\lambda|\ 
}{\lambda\text{ auch komplex }}\)} \\
    6 & \(k  \lambda^{k}  \sigma\lbrack k\rbrack\) & \multicolumn{2}{l}{%
\(\frac{z  \lambda}{(z - \lambda)^{2}}\)} & \multicolumn{2}{>{\raggedright\arraybackslash}p{(\linewidth - 10\tabcolsep) * \real{0.3112} + 2\tabcolsep}@{}}{%
\({|z| > |\lambda|\ 
}{\lambda\text{ auch komplex }}\)} \\
    7 & \(k^{2}  \lambda^{k}  \sigma\lbrack k\rbrack\) & \multicolumn{2}{l}{%
\(\frac{z  \lambda  (z + \lambda)}{(z - \lambda)^{3}}\)} & \multicolumn{2}{>{\raggedright\arraybackslash}p{(\linewidth - 10\tabcolsep) * \real{0.3112} + 2\tabcolsep}@{}}{%
\({|z| > |\lambda|\ 
}{\lambda\text{ auch komplex }}\)} \\
    8 & \(k^{3}  \lambda^{k}  \sigma\lbrack k\rbrack\) & \multicolumn{2}{l}{%
\(\frac{z  \lambda  \left( z^{2} + 4  \lambda  z + \lambda^{2} \right)}{(z - \lambda)^{4}}\)} & \multicolumn{2}{l}{%
\({|z| > |\lambda|\ 
}{\lambda\text{ auch komplex }}\)} \\
    9 & \multicolumn{2}{l}{%
\(\left( \begin{array}{r}
k - 1 \\
    n - 1
\end{array} \right)  \lambda^{k - n}  \sigma\lbrack k - n\rbrack\)} & \multicolumn{2}{l}{%
\(\frac{1}{(z - \lambda)^{n}}\)} & \({|z| > |\lambda|\ 
}{\lambda\text{ auch komplex }}\) \\
    10 & \multicolumn{2}{l}{%
\(\sin\left( k  \Omega_{0} \right)  \sigma\lbrack k\rbrack\)} & \multicolumn{2}{>{\raggedright\arraybackslash}p{(\linewidth - 10\tabcolsep) * \real{0.3131} + 2\tabcolsep}}{%
\(\frac{z  \sin\left( \Omega_{0} \right)}{z^{2} - 2  z  \cos\left( \Omega_{0} \right) + 1}\)} & \(|z| > 1\) \\
    11 & \multicolumn{2}{l}{%
\(\cos\left( k  \Omega_{0} \right)  \sigma\lbrack k\rbrack\)} & \multicolumn{2}{>{\raggedright\arraybackslash}p{(\linewidth - 10\tabcolsep) * \real{0.3131} + 2\tabcolsep}}{%
\(\frac{z  \left( z - \cos\left( \Omega_{0} \right) \right)}{z^{2} - 2  z  \cos\left( \Omega_{0} \right) + 1}\)} & \(|z| > 1\) \\
    12 & \multicolumn{2}{l}{%
\(\frac{1}{k!}  \sigma\lbrack k\rbrack\)} & \multicolumn{2}{>{\raggedright\arraybackslash}p{(\linewidth - 10\tabcolsep) * \real{0.3131} + 2\tabcolsep}}{%
\(e^{\frac{1}{z}}\)} & \(|z| > 1\) \\
    \bottomrule
  \end{tabular}
\end{table}

\section{Rücktransformation}\label{sec:5.3}

Für den Einsatz der z-Transformation bei der Lösung linearer Differenzengleichungen ist eine inverse z-Transformation für die Rücktransformation erforderlich. Sie lässt sich zunächst als mathematische Umkehrformel angeben. Alternativ kann die Rücktransformation mit der Partialbruchzerlegung auf bekannte Korrespondenzen zurückgeführt werden. Als weitere Variante kann ein Potenzreihenansatz für die Rücktransformation gewählt werden.

\subsection{Definition der inversen z-Transformation}\label{sec:definition-der-inversen-z-transformation}

Die Berechnung einer Folge \(x[k]\) aus ihrer z-Transformierten ergibt sich ohne Beweis über ein Umlaufintegral {[}Foel03{]}
\begin{align}
x[ k] &= \frac{1}{2  \pi  j}  \oint_{}^{}{X(z)  z^{k - 1}\ dz} \label{eq:5.82}
\end{align}
Gleichung \eqref{eq:5.82} zur inversen z-Transformation wird an dieser Stelle genutzt, um die Rechenregel zur Multiplikation von Folgen zu beweisen. Die z-Transformierte ist definiert als
\begin{align}
Z\left\{ x[ k]  y[ k] \right\} &= \sum\limits_{k = 0}^{\infty}{x[ k]  y[ k]  z^{- k}} \label{eq:5.83}
\end{align}
Wird die Folge \(x[k]\) als Umlaufintegral ausgedrückt, ergibt sich
\begin{align}
{\sum\limits_{k = 0}^{\infty}{x[ k]  y[ k]  z^{- k}} = \frac{1}{2  \pi  j}  \sum\limits_{k = 0}^{\infty}{\left( \oint_{}^{}{X(v)  v^{k - 1}\ dv} \right)  y[ k]  z^{- k}}
}{		 = \frac{1}{2  \pi  j}  \sum\limits_{k = 0}^{\infty}{y[ k]  \oint_{}^{}{X(v)  v^{- 1}  \left( \frac{z}{v} \right)^{- k}\ dv}}} \label{eq:5.84}
\end{align}
Mit einigen Umformungen kann die Modulationsregel angewendet werden
\begin{align}
{\sum\limits _{k=0}^{\infty }x\left[k\right] y\left[k\right] z^{-k}  } &=  {\frac{1}{2 \pi  j}  \sum\limits _{k=0}^{\infty }y\left[k\right] \oint X\left(v\right) v^{-1}  \left(\frac{z}{v} \right)^{-k} {\; }dv  } \\
{} &=  {\frac{1}{2 \pi  j}  \oint \sum\limits _{k=0}^{\infty }y\left[k\right] \left(\frac{z}{v} \right)^{-k}   X\left(v\right) v^{-1} {\; }dv } \\
{} &=  {\frac{1}{2 \pi  j}  \oint Y\left(\frac{z}{v} \right) X\left(v\right) v^{-1} {\; }dv }
\label{eq:5.85}
\end{align}
Damit ist die Multiplikationsregel hergeleitet. Wie bereits bei der inversen Laplace-Transformation wird die Definition der Rücktransformation aber in der Praxis nicht verwendet. Stattdessen wird die Funktion \(X(z)\) so umgeformt, dass bekannte Korrespondenzen genutzt werden können. Dieser Ansatz führt zur Rücktransformation über eine Partialbruchzerlegung oder über einen Potenzreihenansatz.

\subsection{Rücktransformation über Partialbruchzerlegung}\label{sec:ruxfccktransformation-uxfcber-partialbruchzerlegung}

In den bisher behandelten Beispielen und Rechenregeln sind immer gebrochen rationale Funktionen der Form
\begin{align}
X(z) &= \frac{B(z)}{A(z)} = \frac{\sum\limits_{m = 0}^{M}{b_{m}  z^{m}}}{\sum\limits_{n = 0}^{N}{a_{n}  z^{n}}} = \frac{b_{M}}{a_{N}}  \frac{\prod\limits_{m = 1}^{M}\left( z - \beta_{m} \right)}{\prod\limits_{n = 1}^{N}\left( z - \alpha_{n} \right)} \label{eq:5.86}
\end{align}
entstanden. Da es sich bei \(x[k]\) um ein kausales Signal handelt, ist der Zählergrad M maximal so groß wie der Nennergrad N (M ≤ N). Die Koeffizienten \(a_{n}\) und \(b_{m}\) sind reelle Koeffizienten. Damit sind die Polstellen \(α_{n}\) und Nullstellen \(β_{m}\) von \(X(z)\) entweder reell oder konjugiert komplex zueinander.

Diese gebrochen rationale Funktionen \(X(z)\) lassen sich nur in seltenen Fällen direkt über eine bekannte Korrespondenz zurücktransformieren. Im Allgemeinen ist eine Zerlegung der Funktion mit der Partialbruchzerlegung notwendig. Nach der Zerlegung liegen einzelne Partialbrüche vor, die auf bekannte Korrespondenzen zurückgeführt werden können.

\subsubsection{Vorbereitung der Partialbruchzerlegung falls Zählergrad M gleich Nennergrad N}\label{sec:vorbereitung-der-partialbruchzerlegung-falls-zuxe4hlergrad-m-gleich-nennergrad-n}

Falls der Zählergrad M und der Nennergrad N gleich groß sind, muss vor der Partialbruchzerlegung eine Polynomdivision durchgeführt werden. Dadurch entsteht ein konstanter Summand
\begin{align}
X_{0} &= \frac{b_{M}}{a_{N}} \label{eq:5.87}
\end{align}
Da die inverse z-Transformierte zu einer Konstanten die Impulsfolge ist, entspricht diesem Summanden im Zeitbereich der Folgenwert für \(k = 0\)
\begin{align}
x[ k] &= \frac{b_{M}}{a_{N}}  \delta[ k] \label{eq:5.88}
\end{align}
Das folgende Beispiel verdeutlicht diesen Zusammenhang.

\begin{example}{Zählergrad gleich Nennergrad}
\label{example:beispiel-zuxe4hlergrad-gleich-nennergrad}

Die z-Transformierte \(X(z)\) soll in den Zeitbereich zurück transformiert werden. Da der Zählergrad genauso groß ist wie der Nennergrad, wird eine Polynomdivision durchgeführt. Um auf einen Ausdruck der Korrespondenztafel zu kommen, muss der Restbruch mit z erweitert werden.
\begin{align*}
X(z) &= \frac{2  z - 3}{z - 3} = 2 + \frac{3}{z - 3} = 2 + \frac{3  z}{z - 3}  z^{- 1} 
\end{align*}
Damit kann die Folge im Zeitbereich aus der Korrespondenztafel und der Verschiebungsregel bestimmt werden zu
\begin{align*}
x[ k] &= 2  \delta[ k] + 3  3^{k - 1}  \sigma[ k - 1] 
\end{align*}
\end{example}

\textbf{\hfill\break
}

\subsubsection{Partialbruchzerlegung für einfache Pole bei z ≠ 0}\label{sec:partialbruchzerlegung-fuxfcr-einfache-pole-bei-z-0}

Besitzt die z-Transformierte \(X(z)\) nur einfache Pole, kann Sie mit Hilfe der Partialbruchzerlegung dargestellt werden als
\begin{align}
X(z) &= \frac{1}{a_{N}}  \frac{\sum\limits_{M = 0}^{M}{b_{m}  z^{m}}}{\prod\limits_{n = 1}^{N}\left( z - \alpha_{n} \right)} = \sum\limits_{n = 1}^{N}\frac{A_{n}}{z - \alpha_{n}} \label{eq:5.91}
\end{align}
Die Koeffizienten \(A_{n}\) der einzelnen Partialbrüche können wie bei der Laplace-Transformation auf unterschiedliche Arten berechnet werden:

\begin{itemize}
\item
  Ausmultiplizieren\\
  Die Gleichung wird mit den Polstellen multipliziert. Anschließend werden die Polstellen eingesetzt, und es ergibt sich ein Gleichungssystem für die Koeffizienten \(A_{n}\).
\item
  Residuensatz\\
  Die einzelnen Koeffizienten werden über den Residuensatz berechnet, der grundsätzlich auf dem gleichen Verfahren basiert
\end{itemize}
\begin{align}
A_{n} &= \left. \ \left( X(z)  \left( z - \alpha_{n} \right) \right) \right|_{z = \alpha_{n}} \label{eq:5.92}
\end{align}
Jeder einzelne Partialbruch hat die Form
\begin{align}
X_{n}(z) &= \frac{A_{n}}{z - \alpha_{n}} = A_{n}  \frac{z}{z - \alpha_{n}}  z^{- 1} \label{eq:5.93}
\end{align}
Im Zeitbereich ergibt sich damit für jeden Partialbruch eine kausale Exponentialfolge, die um einen Folgenindex nach rechts verschoben ist:
\begin{align}
x_{n}[ k] &= Z^{- 1}\left\{ \frac{A_{n}}{z - \alpha_{n}} \right\} = A_{n}  \alpha_{n}^{k - 1}  \sigma[ k - 1] \label{eq:5.94}
\end{align}
Die Summe der einzelnen Partialbrüche entspricht deshalb im Zeitbereich der Folge
\begin{align}
Z^{- 1}\left\{ \sum\limits_{n = 1}^{N}\frac{A_{n}}{z - \alpha_{n}} \right\} &= \sum\limits_{n = 1}^{N}{A_{n}  \alpha_{n}^{k - 1}  \sigma[ k - 1]} \label{eq:5.95}
\end{align}
\begin{example}{Partialbruchzerlegung für einfache Pole bei z ≠ 0}
\label{example:beispiel-partialbruchzerlegung-fuxfcr-einfache-pole-bei-z-0}

Die z-Transformierte \(X(z)\) soll in den Zeitbereich zurücktransformiert werden. Ihr Zählergrad ist kleiner als der Nennergrad, und sie hat zwei einfache Pole.
\begin{align*}
X(z) &= \frac{z}{z^{2} + 3  z + 2} = \frac{z}{(z + 1)  (z + 2)} = \frac{A_{1}}{z + 1} + \frac{A_{2}}{z + 2} 
\end{align*}
Ausmultiplizieren der Gleichung führt zu
\begin{align*}
z &= A_{1}  (z + 2) + A_{2}  (z + 1) 
\end{align*}
Einsetzen der Polstellen \(α_{1}\) = - 1 und \(α_{2}\) = - 2 führt zu
\begin{align*}
A_{1} &= - 1 
\end{align*}
und
\begin{align*}
A_{2} &= 2 
\end{align*}
Alternativ hätte der Residuensatz ergeben
\begin{align*}
A_{1} &= \left. \ \left( \frac{z}{(z + 1)  (z + 2)}  (z + 1) \right) \right|_{z = - 1} = \left. \ \left( \frac{z}{z + 2} \right) \right|_{z = - 1} = \frac{- 1}{- 1 + 2} = - 1 
\end{align*}
und
\begin{align*}
A_{2} &= \left. \ \left( \frac{z}{(z + 1)  (z + 2)}  (z + 2) \right) \right|_{z = - 2} = \left. \ \left( \frac{z}{z + 1} \right) \right|_{z = - 2} = \frac{- 2}{- 2 + 1} = 2 
\end{align*}
Sind die Koeffizienten der Partialbrüche bestimmt, kann die z-Transformierte mit den bekannten Korrespondenzen in den Zeitbereich zurücktransformiert werden.
\begin{align*}
X(z) &= \frac{- 1}{z + 1} + \frac{2}{z + 2} = \frac{- z}{z + 1}  z^{- 1} + \frac{2  z}{z + 2}  z^{- 1} 
\end{align*}
Es ergibt sich die Folge
\begin{align*}
x[ k] &= - 1  ( - 1)^{k - 1}  \sigma[ k - 1] + 2  ( - 2)^{k - 1}  \sigma[ k - 1] 
\end{align*}
\end{example}

\subsubsection{Partialbruchzerlegung für mehrfache Pole bei z ≠ 0}\label{sec:partialbruchzerlegung-fuxfcr-mehrfache-pole-bei-z-0}

Liegt ein P-facher Pol an der Stelle α vor, muss der Teil der z-Transformierten dargestellt werden als
\begin{align}
X(z) &= \frac{B(z)}{(z - \alpha)^{p}} = \sum\limits_{n = 1}^{P}\frac{A_{n}}{(z - \alpha)^{n}} \label{eq:5.104}
\end{align}
Die Koeffizienten \(A_{n}\) der einzelnen Partialbrüche können wieder auf unterschiedliche Arten berechnet werden:

\begin{itemize}
\item
  Ausmultiplizieren\\
  Die Gleichung wird mit den Polstellen multipliziert. Anschließend wird die Polstelle und P - 1 weitere Werte für z eingesetzt, und es ergibt sich ein Gleichungssystem für die\\
  Koeffizienten \(A_{n}\).
\item
  Residuensatz\\
  Die einzelnen Koeffizienten werden über den Residuensatz berechnet
\end{itemize}
\begin{align}
A_{n} &= \frac{1}{(P - n)!}  \frac{d^{P - n}}{dz^{P - n}}\left. \ \left( X(z)  (z - \alpha)^{P} \right) \right|_{z = \alpha} \label{eq:5.105}
\end{align}
Die Rücktransformation kann mit Korrespondenz 9 angegeben werden zu
\begin{align}
x[ k] &= \sum\limits_{n = 1}^{P}{A_{n}  \left( \begin{array}{r}
k - 1 \\
n - 1
\end{array} \right)  \alpha^{k - n}  \sigma[ k - n]} \label{eq:5.106}
\end{align}
\begin{example}{Partialbruchzerlegung für mehrfache Pole bei z ≠ 0}
\label{example:beispiel-partialbruchzerlegung-fuxfcr-mehrfache-pole-bei-z-0}

Die z-Transformierte \(X(z)\) soll in den Zeitbereich zurücktransformiert werden. Ihr Zählergrad ist kleiner als der Nennergrad und sie hat einen doppelten Pol an der Stelle \(α = 0.5\). Damit lautet der Ansatz für die Partialbruchzerlegung
\begin{align*}
X(z) &= \frac{z}{(z - 0.5)^{2}} = \frac{A_{1}}{z - 0.5} + \frac{A_{2}}{(z - 0.5)^{2}} 
\end{align*}
Ausmultiplizieren führt zu der Gleichung
\begin{align*}
z &= A_{1}  (z - 0.5) + A_{2} 
\end{align*}
Einsetzen der Zahlenwerte \(z = 0.5\) und \(z = 0\) ergibt
\begin{align*}
A_{2} &= 0.5 
\end{align*}
und
\begin{align*}
A_{1} &= 2  A_{2} = 1 
\end{align*}
Alternativ könnte der Residuensatz verwendet werden:
\begin{align*}
A_{1} &= \frac{d}{dz}\left. \ \left( X(z)  (z - 0.5)^{2} \right) \right|_{z = 0.5} = \left. \ \frac{dz}{dz} \right|_{z = 0.5} = 1 
\end{align*}
und
\begin{align*}
A_{2} &= 1  \left. \ X(z)  (z - 0.5)^{2} \right|_{z = z_{\infty}} = \left. \ z \right|_{z = 0.5} = 0.5 
\end{align*}
Die Funktion \(X(z)\) kann damit in folgende Partialbrüche aufgeteilt werden:
\begin{align*}
X(z) &= \frac{z}{(z - 0.5)^{2}} = \frac{1}{z - 0.5} + \frac{0.5}{(z - 0.5)^{2}} = \frac{z}{z - 0.5}  z^{- 1} + \frac{0.5  z}{(z - 0.5)^{2}}  z^{- 1} 
\end{align*}
Zur Rücktransformation werden die Korrespondenzen 5 und 6 sowie der Verschiebungssatz verwendet.
\begin{align*}
x[ k] &= 0.5^{k - 1}  \sigma[ k - 1] + (k - 1)  0.5^{k - 1}  \sigma[ k - 1] 
\end{align*}
Alternativ ist eine Rücktransformation mit Korrespondenz 9 möglich.
\begin{align*}
{x\left[k\right]} &=  {\left(\begin{array}{c} {k-1} \\
{0} \end{array}\right) 0.5^{k-1-1+1}  \sigma \left[k-1-1+1\right]+0.5 \left(\begin{array}{c} {k-1} \\
{1} \end{array}\right) 0.5^{k-1-2+1}  \sigma \left[k-1-2+1\right]} \\
{} &=  {0.5^{k-1}  \sigma \left[k-1\right]+(k-1) 0.5^{k-1}  \sigma \left[k-2\right]}
\end{align*}
Da der zweite Summand für \(k = 1\) zu null wird, sind beide Lösungen identisch.

\end{example}

\subsubsection{Sonderfall N-facher Pol bei \(z = 0\)}\label{sec:sonderfall-n-facher-pol-bei-z-0}

Wenn nur Pole an der Stelle \(z = 0\) vorliegen, so sind alle Koeffizienten \(a_{n} = 0\) bis auf den Koeffizienten \(a_{N}\). Damit gilt
\begin{align}
X(z) &= \frac{\sum\limits_{m = 0}^{M}{b_{m}  z^{m}}}{a_{N}  z^{N}} = \sum\limits_{m = 0}^{M}{\frac{b_{m}}{a_{N}}  z^{m - N}} = \frac{b_{0}}{a_{N}}  z^{- N} + \frac{b_{1}}{a_{N}}  z^{1 - N} + ... + \frac{b_{M}}{a_{N}}  z^{M - N} \label{eq:5.116}
\end{align}
Da bei kausalen Signalen alle Exponenten der komplexen Variablen z immer M - N \textless{} 0 sind, kann der Verschiebungssatz der z-Transformation für eine Verschiebung nach rechts angewendet werden. Die Folge \(x[k]\) im Zeitbereich ist damit
\begin{align}
x[ k] &= \frac{b_{0}}{a_{N}}  \delta[ k - N] + \frac{b_{1}}{a_{N}}  \delta[ k - N + 1] + ... + \frac{b_{M}}{a_{N}}  \delta[ k - N + M] \label{eq:5.117}
\end{align}
\begin{example}{N-facher Pol an der Stelle \(z = 0\)}
\label{example:beispiel-n-facher-pol-an-der-stelle-z-0}

Die z-Transformierte \(X(z)\) soll in den Zeitbereich zurücktransformiert werden. Sie hat einen 4-fachen Pol an der Stelle \(z = 0\).
\begin{align*}
X(z) &= \frac{z^{4} + 2  z^{3} + 3  z^{2} + 2  z + 1}{9  z^{4}} = \frac{1}{9}  \left( 1 + 2  z^{- 1} + 3  z^{- 2} + 2  z^{- 3} + 1  z^{- 4} \right) 
\end{align*}
Mit dem Verschiebungssatz der z-Transformation ergibt sich im Zeitbereich die Folge
\begin{align*}
x[ k] &= \frac{1}{9}  \delta[ k] + \frac{2}{9}  \delta[ k - 1] + \frac{3}{9}  \delta[ k - 2] + \frac{2}{9}  \delta[ k - 3] + \frac{1}{9}  \delta[ k - 4] 
\end{align*}
\end{example}

\subsubsection{Zusammenfassung der Ansätze für die Partialbruchzerlegung}\label{sec:zusammenfassung-der-ansuxe4tze-fuxfcr-die-partialbruchzerlegung}

Tabelle \ref{tab:5.6} fasst die Ansätze für die Partialbruchzerlegung zusammen. Dabei wird von einer\\
z-Transformierten
\begin{align}
X(z) &= \frac{B(z)}{A(z)} = \frac{\sum\limits_{m = 0}^{M}{b_{m}  z^{m}}}{\sum\limits_{n = 0}^{N}{a_{n}  z^{n}}} = \frac{b_{M}}{a_{N}}  \frac{\prod\limits_{m = 1}^{M}\left( z - \beta_{m} \right)}{\prod\limits_{n = 1}^{N}\left( z - \alpha_{n} \right)} \label{eq:5.120}
\end{align}
ausgegangen, bei der der Zählergrad M kleiner als der Nennergrad N ist. Die Koeffizienten \(a_{n}\) und \(b_{m}\) sind reelle Koeffizienten. Die Nullstellen \(β_{m}\) und die Pole \(α_{n}\) sind nicht gleich.

\begin{table}
  \centering
  \caption{Ansätze für die Partialbruchzerlegung}
  \label{tab:5.6}
  \begin{tabular}{ll}
    \toprule
    Pollage & Ansatz Partialbruchzerlegung \\
    \midrule
    Einfache reelle oder komplexe Pole

\(\alpha_{n}\) & \(X(z) = \sum\limits_{n = 1}^{N}\frac{A_{n}}{z - \alpha_{n}}\) \\
    p-facher reeller oder komplexer Pol

\(\alpha_{n}\) & \(X(z) = \sum\limits_{n = 1}^{P}\frac{A_{n}}{(z - \alpha)^{n}}\) \\
    Ausschließlich Polstellen an der Stelle \(z = 0\) & \(X(z) = \sum\limits_{m = 0}^{M}{\frac{b_{m}}{a_{N}}  z^{m - N}}\) \\
    \bottomrule
  \end{tabular}
\end{table}

\subsection{Rücktransformation über Reihenentwicklung}\label{sec:ruxfccktransformation-uxfcber-reihenentwicklung}

Funktionen können über Potenzreihen beschrieben werden. Für gebrochen rationale Funktionen ergeben sich die Potenzreihen aus einer Polynomdivision von Zähler- und Nennerpolynom. Diese Division ist im allgemeinen Fall jedoch nicht endlich. Es kann ein Polynom mit unendlich vielen Summanden entstehen, die mit der Verschiebungsregel rücktransformiert werden können.

Die Darstellung von Funktionen \(X(z)\) über sogenannte Laurent-Reihen sowie die entsprechende Rücktransformation sind in {[}Oppe04{]} beschrieben.

\begin{example}{Inverse z-Transformation über Reihenentwicklung}
\label{example:beispiel-inverse-z-transformation-uxfcber-reihenentwicklung}

Die Funktion \(X(z)\) soll mit Hilfe einer Reihenentwicklung in den Zeitbereich transformiert werden. Durch Polynomdivision werden die Glieder der Potenzreihe berechnet zu
\begin{align*}
X(z) &= \frac{z}{z + 0.5} = 1 + 0.5  z^{- 1} + 0.25  z^{- 2} + 0.125  z^{- 3} + 0.0625  z^{- 4} + ... 
\end{align*}
Rücktransformation ergibt mit Anwendung der Verschiebungsregel
\begin{align*}
x[ k] &= \delta[ k] + 0.5  \delta[ k - 1] + 0.25  \delta[ k - 2] + 0.125  \delta[ k - 3] + 0.0625  \delta[ k - 4] + ... 
\end{align*}
An diesem Beispiel zeigt sich, dass die Potenzreihe nicht endlich ist. Für eine exakte Lösung von \(x[k]\) müssten unendlich viele Glieder berechnet werden.

\end{example}

\protect\phantomsection\label{_Toc177369030}{}
